\medterm{YAG Capsulotomy} A laser procedure that makes a small opening in the cloudy membrane behind an artificial lens after cataract surgery. It can restore vision when this clouding causes blurring or glare and is usually done without another incision.

\medterm{YAG Laser Capsulotomy} Alternate name; see \textit{YAG Capsulotomy}.

\medterm{YAG Laser Surgery} Treatment using a focused laser to make a precise opening or alter tissue without a conventional surgical incision. It is used mainly for selected eye conditions, and risks depend on the part of the eye treated.

\medterm{Yawn} An involuntary behavior consisting of a deep inhalation followed by a slower exhalation, often with stretching of the jaw and facial muscles. Yawning is common with fatigue but excessive or unusual yawning can occur with some medications or neurologic conditions.

\medterm{Yawning} Alternate name; see \textit{Yawn}.

\medterm{Yaws} Yaws is a chronic infection caused by Treponema pallidum subsp. pertenue and spread mainly through close skin contact. It can cause an initial papilloma or ulcer and later destructive lesions of skin, bone, and cartilage; antibiotics can cure the infection.

\medterm{Yaws Papilloma} The primary clinical cutaneous lesion of early yaws (\textit{Treponema pallidum} subsp. \textit{pertenue} infection), also termed a mother yaw or frambesioma. It presents as an initial painless, raised, yellow-crusted, raspberry-like papule or ulcer at the site of cutaneous inoculation.

\synonyms
Mother Yaw; Frambesioma; Primary Yaws Lesion

\medterm{Y-Box Binding Protein} A protein that attaches to DNA and RNA and helps control which genes are active and how proteins are made. It also helps cells respond to stress and repair damaged DNA.

\synonyms
YB-1; Y-Box Transcription Factor; Nuclease-Sensitive Element-Binding Protein 1

\medterm{Y Chromosome} One of the two human sex chromosomes. It usually carries genes involved in testis development and is passed by a biological father to a child who inherits it; chromosome patterns and body development can vary.

\medterm{Y-Chromosome Microdeletion} A small missing section of the Y chromosome that can interfere with sperm production and cause very low or absent sperm counts.

\synonyms
AZF Microdeletion

\medterm{Yeast} A type of fungus made of single cells. Some yeasts normally live on or in the body, while others can cause infections of the mouth, skin, vagina, lungs, blood, or other organs.

\medterm{Yeast Infection} An infection caused by Candida yeast that may affect the mouth, skin, vagina, penis, or other tissues. Symptoms depend on the site and may include itching, soreness, redness, white plaques, discharge, or rash. Candida may also be present without causing disease, so the significance of a finding depends on the site, symptoms, and evidence of tissue invasion.

\medterm{Yeast Infection In Women And Men} Genital Candida infection may affect the vulva or vagina in women and the glans penis or foreskin in men. Symptoms can include itching, soreness, redness, discharge, or rash.

\medterm{Yeast Present} A laboratory report showing yeast in a specimen. This may mean harmless colonization, contamination, or infection; the meaning depends on the symptoms, specimen, and type of yeast.

\medterm{Yeast Vulvitis} Inflammation of the vulva associated with a Candida yeast infection. It may cause itching, burning, redness, swelling, and soreness; other conditions can cause similar symptoms.

\medterm{Yellow Cartilage} A flexible type of cartilage containing many elastic fibers. It is found in the outer ear, the epiglottis, parts of the voice box, and the tube connecting the middle ear to the throat.

\synonyms
Elastic Cartilage; Cartilago Elastica

\medterm{Yellow Fever} Yellow fever is an acute mosquito-borne infection caused by yellow fever virus. Fever, headache, muscle aches, nausea, and vomiting may be followed by jaundice, bleeding, shock, and organ failure; vaccination is the main preventive measure.

\medterm{Yellow Fever Vaccine} A vaccine made from a weakened yellow fever virus. It gives long-lasting protection, but rare serious reactions can affect the brain or internal organs.

\synonyms
17D Vaccine

\medterm{Yellow Fever Virus} A flavivirus transmitted mainly by infected mosquitoes that can cause yellow fever, ranging from a mild febrile illness to jaundice, bleeding, shock, and organ failure.

\medterm{Yellow Ligament} A series of elastic bands connecting the back parts of neighboring bones in the spine. They help support the spine and surround the spinal canal.

\synonyms
Ligamentum Flavum; Flavum Ligament

\medterm{Yellow Marrow} Bone marrow containing mostly fat instead of blood-forming cells. It is common in the hollow centers of adult long bones and can change back toward blood-forming marrow when the body needs more blood cells.

\synonyms
Fatty Bone Marrow; Medulla Ossium Flava

\medterm{Yellow Nail Syndrome} Alternate name; see \textit{Yellow-Nail Syndrome}.

\medterm{Yellow-Nail Syndrome} A rare disorder causing yellow, thickened, slowly growing nails together with swelling from lymph fluid, usually in the legs, and lung problems such as repeated infections or fluid around the lungs. The cause is not fully understood.

\medterm{Yellow Nail Syndrome Triad} The three main features of yellow nail syndrome: slow-growing yellow or thickened nails, swelling from lymph-fluid buildup, and long-term breathing or lung problems.

\synonyms
YNS Triad; Lymphedema-Pleural Effusion-Yellow Nails

\medterm{Yellow Spot} The macula, a small yellowish area near the center of the retina at the back of the eye. It provides the sharp central vision used for reading and seeing fine detail.

\synonyms
Macula Lutea; Anatomical Macula

\medterm{Yellow Tongue} Yellow discoloration or coating of the tongue that may result from oral hygiene, smoking, dry mouth, diet, medicines, or an underlying illness.

\medterm{Yergason Test} A shoulder examination in which a person turns the palm upward and rotates the arm outward against resistance. Pain or a snapping sensation near the front of the shoulder may suggest a problem involving the biceps tendon, but the test is not diagnostic by itself.

\medterm{Yersinia} A group of bacteria that can cause intestinal illness, swollen lymph nodes in the abdomen, or plague. Different species spread through contaminated food or water, infected animals, or fleas.

\medterm{Yersinia Enterocolitica} A bacterium that can cause diarrhea, fever, and abdominal pain, sometimes mimicking appendicitis. People usually acquire it from contaminated food, water, or contact with infected animals.

\medterm{Yersinia Enterocolitica Pseudoappendicitis} A Yersinia infection that causes pain in the lower right abdomen and swollen lymph nodes near the intestine, making it resemble appendicitis. It can also cause fever and diarrhea, often after contaminated food such as undercooked pork.

\synonyms
Yersinia Mesenteric Adenitis; Pseudoappendiceal Syndrome

\medterm{Yersinia Infection} Infection caused by Yersinia bacteria, usually producing diarrhea, abdominal pain, fever, or vomiting. Some cases resemble appendicitis; bloody diarrhea, dehydration, or severe or persistent illness can occur.

\medterm{Yersinia Outer Proteins} Proteins made by disease-causing Yersinia bacteria that are injected into human cells. They weaken the immune response and help the bacteria avoid being destroyed.

\synonyms
Yops; Yersinia Effector Proteins; Yop Virulence Factors

\medterm{Yersinia Pestis} The bacterium that causes plague, a severe infection usually spread by infected fleas or animals. It may cause painful swollen lymph nodes, bloodstream infection, or lung disease and can be life-threatening.

\medterm{Yersinia Pestis Pneumonic Plague} A severe form of plague affecting the lungs. It can spread through breathing in infected droplets or when plague spreads from another part of the body. It may cause high fever, breathing difficulty, chest pain, cough, shock, and rapidly worsening lung failure.

\synonyms
Pneumonic Plague; Pulmonary Yersiniosis

\medterm{Yersinia Pseudotuberculosis} A bacterium usually acquired from contaminated food or water. It can cause fever, diarrhea, and pain from swollen lymph nodes in the abdomen, sometimes resembling appendicitis.

\medterm{Yersinia Reactive Arthritis} Joint inflammation that develops after a Yersinia intestinal infection, usually one to four weeks later. It often affects a few joints, especially in the legs, without bacteria being present inside the joints.

\medterm{Yersiniosis} Infection caused by Yersinia bacteria, most often Y. enterocolitica or Y. pseudotuberculosis, producing diarrhea, abdominal pain, fever, or pseudoappendicitis.

\medterm{Yew Poisoning} Poisoning after eating parts of the yew plant, especially its leaves or seeds. The toxins can cause vomiting, dizziness, seizures, dangerous heart rhythms, collapse, and death.

\medterm{Y-Linked Gene} A gene located specifically on the non-recombining region of the Y chromosome (such as the \textit{SRY} gene, \textit{ZFY}, or \textit{AZF} azoospermia factor region). These genes are transmitted strictly from biological fathers to sons and are essential for testicular differentiation, gonadal development, and male spermatogenesis.

\synonyms
Holandric Gene; Y-Chromosomal Gene

\medterm{Y-Linked Inheritance} Inheritance of a condition caused by a gene on the Y chromosome. A biological father passes a Y-linked gene to sons who inherit his Y chromosome, but not to daughters who inherit his X chromosome.

\medterm{Y-Maze Test} A test used mainly in laboratory animals to study movement, exploration, and short-term memory. The animal moves through three arms of a Y-shaped maze, and the pattern of movement is recorded.

\synonyms
Y-Maze Spontaneous Alternation; Y-Maze Behavioral Paradigm

\medterm{Yolk Sac} A temporary structure that supports an early embryo before the placenta is fully working. It helps form the first blood cells and part of the early intestine and can be seen on early pregnancy ultrasound.

\medterm{Yolk Sac Carcinoma} An aggressive germ-cell cancer that often produces high levels of alpha-fetoprotein (AFP). AFP can support the diagnosis and help track treatment, but it is interpreted with examination, imaging, and tissue testing.

\synonyms
Endodermal Sinus Tumor; Yolk Sac Germ Cell Tumor

\medterm{Yolk Sac Tumor} A cancer of germ cells that often occurs in the ovary or testicle in children or young adults and produces alpha-fetoprotein. It can also occur outside the reproductive organs.

\medterm{Yolk-Sac Tumor} Alternate name; see \textit{Yolk Sac Tumor}.

\medterm{Yolk Sac Tumor Schiller-Duval Bodies} A characteristic microscope finding in yolk sac tumors. It looks somewhat like a tiny kidney filter, with tumor cells arranged around a small blood vessel. The finding helps confirm the diagnosis when considered with other test results.

\synonyms
Schiller-Duval Bodies; Glomeruloid Endodermal Bodies

\medterm{Yolk Stalk} A temporary connection between the early yolk sac and the developing intestine. It normally disappears before birth; a remaining portion can form an abnormal intestinal pouch or other problem.

\medterm{Young-Burgess Classification} A system used to describe pelvic-ring fractures according to how the injury happened and how stable the pelvis is. Groups include side-to-side compression, front-to-back compression, vertical separation, and mixed injuries.

\synonyms
Young-Burgess Pelvic Fracture Classification; Burgess Pelvic Injury Scale

\medterm{Young Syndrome} A rare condition, usually affecting males, that combines blockage of sperm transport with repeated sinus and lung infections. It may also cause bronchiectasis, a long-term widening and damage of the airways.

\synonyms
Sinopulmonary-Infertility Syndrome; Barry-Perkins-Young Syndrome

\medterm{Y-Plate Osteosynthesis} Surgical repair of a complex broken bone using a metal plate shaped like the letter Y. The plate holds the bone pieces in place while they heal.

\synonyms
Y-Plate Fixation; Y-Shaped Contoured Bone Plate

\medterm{Yttrium-90} A radioactive substance used in selected treatments to deliver radiation to abnormal tissue, such as some liver tumors or inflamed joints.

\medterm{Yttrium-90 Microspheres} Tiny glass or resin beads containing radioactive yttrium-90. They are delivered through blood vessels to selected liver tumors, where they release radiation.

\synonyms
Y-90 Radioembolization Microspheres; TheraSphere; SIR-Spheres

\medterm{Yttrium-90 Radioembolization} A treatment in which tiny beads containing radioactive yttrium-90 are delivered through liver blood vessels to release radiation inside liver tumors.

\synonyms
SIRT; Y-90 SIRT; Y-90 Transarterial Radioembolization

\medterm{Yttrium-Aluminum-Garnet Laser} A laser that uses a crystal to produce a focused beam of light. It is used in eye care for procedures such as opening a cloudy membrane after cataract surgery and making a small opening for some types of glaucoma.

\synonyms
Nd:YAG Laser; Neodymium:YAG Laser

\medterm{Yuzpe Regimen} An older form of emergency contraception using two doses of combined birth-control hormones. It has largely been replaced by emergency contraception containing levonorgestrel or ulipristal.

\medterm{Y-V Plasty} A surgical technique in which a Y-shaped cut is closed as a V. This lengthens or widens nearby tissue and may help release a tight scar or other narrowing.

\synonyms
Y-V Advancement Flap; Y-V Lengthening Procedure
